\chapter{Public Choice: Voting}\label{ch:b3:public-choice-voting}

The previous chapters asked what a government should do if it maximised social
welfare. This one asks what governments do, given that they are chosen by voters,
run by politicians and constrained by rules. The same tools apply: people pursue
their aims under constraints, and the predictions are tested against the record.
Market failures justify a role for government; this chapter and the next measure
how the government that fills the role is likely to behave.

\begin{openingchart}{Central-government expenditure as a share of GDP, 1870--2020:
median and interquartile range of up to 18 economies (years with at least 12), and
the United Kingdom and the United States (percent).}{Jordà, Schularick and Taylor,
Macrohistory Database (expenditure, gdp); our ratios}
\begin{tikzpicture}
\begin{axis}[omtimesten, width=0.72\linewidth, height=0.42\linewidth, xmin=1865,
  xmax=2025, ymin=0, ymax=70, ylabel={Percent of GDP}, legend pos=north west,
  legend style={font=\footnotesize}]
  \addplot[name path=lo, draw=none, forget plot] table[x=year, y=q25] {figdata/bachelor-3/21-public-choice-voting/out/spending.csv};
  \addplot[name path=hi, draw=none, forget plot] table[x=year, y=q75] {figdata/bachelor-3/21-public-choice-voting/out/spending.csv};
  \addplot[omDef!20, area legend] fill between[of=lo and hi];
  \addplot[omDef, thick, mark=none] table[x=year, y=median] {figdata/bachelor-3/21-public-choice-voting/out/spending.csv};
  \addplot[omThm, thick, dashed, mark=none] table[x=year, y=GBR] {figdata/bachelor-3/21-public-choice-voting/out/spending.csv};
  \addplot[black, thick, dotted, mark=none] table[x=year, y=USA] {figdata/bachelor-3/21-public-choice-voting/out/spending.csv};
  \legend{Interquartile range, Median, United Kingdom, United States}
\end{axis}
\end{tikzpicture}
\end{openingchart}

% figdata: spending_points
The median economy's central government spent 6.5~\% of GDP in 1870 and 7.4~\% in 1913,
12.1~\% in 1938, 17.2~\% in 1950, 25.3~\% in 1980 and 31.4~\% in 2020. The wars raised it in steps,
but most of the rise after 1950 came in peacetime, as the vote spread, incomes grew
and governments took on pensions, health care and schooling. The models of this
chapter try to explain that rise and are tested against it.

\begin{recall}
The grades 10--12 volume introduced public choice. The Year 2 volume proved the
median-voter theorem for single-peaked preferences over one dimension and stated
Arrow's theorem. The regression discontinuity design is chapter 4's, and the
capitalisation of taxes into asset prices was met in the Year 1 volume.
\end{recall}

\begin{notation}[Symbols of this chapter]\label{not:b3:public-choice-voting:symbols}
$q$ is the size of a public programme; $q_0$ the \omterm{def:b3:public-choice-voting:agenda-setter}{reversion point}; $t$ a tax rate; $y_m$ and $\bar y$
the median and mean incomes; $P$ the probability that one vote decides an election.
\end{notation}

\section{The median voter and its tests}\label{sec:b3:public-choice-voting:median}

\begin{model}{The median voter}
\assumes Policy has one dimension, voters' preferences over it are single-peaked,
and two parties that care only about winning commit to platforms.
\predicts Both platforms converge to the median voter's ideal point, and policy
moves with the median voter's preferences.
\tested In close elections to the United States House of Representatives,
1946--1994, which party wins is as good as random; the roll-call votes of the
winners do not respond to large changes in their chance of winning, which is
closer to full divergence --- voters choose between fixed platforms --- than to the
convergence the model predicts, full or partial. % ledger: b3pv:lmb
\breaks Politicians care about policy and cannot commit to moderate it; turnout
differs across voters; policy has many dimensions; and incumbents hold an
advantage: in the same elections, a narrow win raised the chance of winning the
next one by about 0.40 to 0.45. % ledger: b3pv:lee
\end{model}

\section{Many dimensions: cycles}\label{sec:b3:public-choice-voting:cycles}

\begin{definition}[Condorcet winner; majority cycle]\label{def:b3:public-choice-voting:cycle}
\index{Condorcet winner}\index{majority cycle}
A \emph{Condorcet winner} is an alternative that beats every other in a pairwise
majority vote. A \emph{majority cycle} is a set of alternatives in which each is
beaten by another, so that none is a Condorcet winner.
\end{definition}

\begin{proposition}[A cycle with three voters]\label{prop:b3:public-choice-voting:cycle}
If three voters rank three alternatives $A \succ B \succ C$, $B \succ C \succ A$ and
$C \succ A \succ B$, majority rule prefers $A$ to $B$, $B$ to $C$ and $C$ to $A$.
\end{proposition}

\begin{proof}
$A$ beats $B$ with the first and third voters; $B$ beats $C$ with the first and
second; $C$ beats $A$ with the second and third. Each pair is decided two to one,
and the three decisions form a cycle.
\end{proof}

The third voter's preferences are not single-peaked on the order $A, B, C$: she
ranks both ends above the middle. When the preferences of every voter are
single-peaked on one dimension, the median voter's ideal beats every alternative and
no cycle can form. Arrow's theorem, recalled above, says that with three or more
alternatives no rule other than dictatorship meets a short list of reasonable
conditions for every profile of preferences. How often do they arise? The listing
draws profiles in three ways: every ranking equally likely, voters who rank
alternatives by distance on a line, and voters who rank them by distance in a plane.

\omcode{code/bachelor-3/21-public-choice-voting/listings/condorcet.py}{1}{32}{The frequency of \omterm{def:b3:public-choice-voting:cycle}{majority cycles} over three alternatives.}\label{lst:b3:public-choice-voting:condorcet}

\omoutput{code/bachelor-3/21-public-choice-voting/outputs/condorcet.txt}{Share of simulated profiles with a cycle.}\label{out:b3:public-choice-voting:condorcet}

With every ranking equally likely, a cycle appears in 12 of the 216 profiles of
three voters, 5.6~\%, and in about 9~\% of large electorates, close to the known limit
of 8.8~\%. On a line, never. In a plane, rarely: voters spread evenly around a centre
leave each alternative's majority anchored near it, and cycles appear in about 1~\%
of profiles or fewer. % output: condorcet
Cycles are a matter of how preferences are spread, not of majority rule alone:
when the electorate has a centre, majority rule has a winner. What the theorem
leaves open is who sets the agenda when it does not, the subject of the section
after next.

\section{Turnout and the uninformed voter}\label{sec:b3:public-choice-voting:turnout}

\begin{definition}[Calculus of voting; rational ignorance]\label{def:b3:public-choice-voting:turnout}
\index{calculus of voting}\index{paradox of voting}\index{rational ignorance}
The \emph{calculus of voting} compares the expected benefit of voting --- the
probability $P$ that one vote decides the outcome times the value of the preferred
outcome --- with its cost; when the benefit is far below the cost while many vote,
the gap is the \emph{paradox of voting}. \emph{Rational ignorance} is a voter's choice
not to acquire costly political information whose expected value to her is small.
\end{definition}

\begin{measure}[The probability of being decisive]
A week before the 2008 United States presidential election, forecasts put the
chance that one vote would decide the national outcome at about 1 in 10 million in
the closest states and 1 in 60 million on average. % ledger: b3pv:gse
At such odds, the private benefit of voting is close to zero unless the voter
values the outcome for millions of others, or values voting itself.
\end{measure}

People vote anyway, in large numbers, which is the paradox; but they spend little
on informing themselves, which is what the calculus predicts. Groups whose members
each have much at stake in one policy --- an industry, a profession --- inform
themselves and organise; the next chapter takes them up.

\section{Vote trading and agenda setting}\label{sec:b3:public-choice-voting:agenda}

\begin{definition}[Logrolling]\label{def:b3:public-choice-voting:logrolling}
\index{logrolling}\index{vote trading}
\emph{Logrolling} is the trading of votes among legislators: each supports
another's project in exchange for support for his own.
\end{definition}

\omterm{def:b3:public-choice-voting:logrolling}{Logrolling} can pass projects whose total cost exceeds their total benefit, when
each benefits a few districts and the cost is shared by all: every member of the
coalition gains from the trade, and the districts outside it pay for projects that
give them nothing (\cref{exo:b3:public-choice-voting:11}).

\begin{definition}[Agenda setter (reversion point)]\label{def:b3:public-choice-voting:agenda-setter}
\index{agenda setter}\index{reversion point}
An \emph{agenda setter} is a body that alone may propose a policy, which voters
accept or reject; the \emph{reversion point} is the policy that holds if they
reject.
\end{definition}

\begin{proposition}[The setter gets more than the median voter wants]\label{prop:b3:public-choice-voting:setter}
Let voters' preferences be single-crossing (a voter with a higher ideal prefers a
larger $q$ to a smaller one whenever a voter with a lower ideal does), the median
voter's be single-peaked with ideal $q^*$, and the \omterm{def:b3:public-choice-voting:agenda-setter}{reversion
point} $q_0 < q^*$. A setter who wants the largest $q$ proposes the largest $q > q^*$ that the median
voter likes as much as $q_0$, and it passes.
\end{proposition}

\begin{proof}
The median voter accepts any $q$ she likes at least as much as $q_0$; by single-peakedness
these form an interval $[q_0, \bar q]$ containing $q^*$, with $\bar q > q^*$ where her utility returns to its
level at $q_0$. By single crossing, every voter whose ideal is above hers prefers $\bar q$ to
$q_0$ whenever she does, and every voter below prefers $q_0$ to a larger proposal whenever she
does; so a majority ranks each proposal against $q_0$ as she does, $\bar q$ passes and no larger
proposal does.
The lower the \omterm{def:b3:public-choice-voting:agenda-setter}{reversion point}, the larger $\bar q$.
\end{proof}

\begin{omchart}{Net benefit of a pool to Brenmoor's median household (two members),
by size in lanes: the ideal is 6 lanes; against a reversion point of 2 lanes, a
setter can win a vote on 13.4.}{Invented data: the weekend problem; our computation}
\begin{tikzpicture}
\begin{axis}[omrecord, width=0.56\linewidth, height=0.36\linewidth, xmin=0, xmax=20,
  ymin=-40, ymax=110, xlabel={Pool size $q$ (lanes)}, ylabel={Dalens a year}]
  \addplot[omDef, thick, mark=none] table[x=q, y=u] {figdata/bachelor-3/21-public-choice-voting/out/brenmoor_u.csv};
  \draw[guide] (axis cs:2,43.18) -- (axis cs:13.42,43.18);
  \fill (axis cs:2,43.18) circle (1.6pt) node[above left, font=\scriptsize] {$q_0$};
  \fill (axis cs:6,95.01) circle (1.6pt) node[above, font=\scriptsize] {ideal};
  \fill (axis cs:13.42,43.18) circle (1.6pt) node[above right, font=\scriptsize] {setter};
\end{axis}
\end{tikzpicture}
\end{omchart}

% figdata: brenmoor_stats, brenmoor_u

\section{The size of government}\label{sec:b3:public-choice-voting:size}

\begin{definition}[Meltzer--Richard model]\label{def:b3:public-choice-voting:meltzer-richard}
\index{Meltzer--Richard model}
The \emph{Meltzer--Richard model} is the model in which a linear income tax finances
an equal transfer to all, taxation has a deadweight cost, and the tax rate is the
one the median voter prefers.
\end{definition}

\begin{proposition}[Redistribution and the median]\label{prop:b3:public-choice-voting:mr}
If the transfer is $t\bar y - \tfrac a2t^2\bar y$, the median voter chooses $t = (1 - y_m/\bar y)/a$ when $y_m < \bar y$, and zero
otherwise: the rate rises with the ratio of mean to median income.
\end{proposition}

\begin{proof}
Her income after tax and transfer is $(1 - t)y_m + t\bar y - \tfrac a2t^2\bar y$. Its derivative in $t$ is $\bar y - y_m -
at\bar y$, zero at $t = (1 - y_m/\bar y)/a$, and the second derivative $-a\bar y$ is negative. If $y_m \ge \bar y$ the derivative
is negative at $t = 0$. As $\bar y/y_m$ rises, $y_m/\bar y$ falls and $t$ rises.
\end{proof}

\begin{model}{Meltzer--Richard}
\assumes Voters differ only in income; the tax is linear, the transfer equal, and
the median voter decides.
\predicts More unequal societies, and societies that extend the vote to poorer
citizens, redistribute more.
\tested Within countries, the franchise prediction holds. When poll taxes and literacy
tests were abolished in U.S. states, turnout rose most among the poor, the pivotal
voter became poorer, and welfare spending rose sharply while other spending did not
(Husted and Kenny 1997, 46 states, 1950--1988). % ledger: b3pv:hk.franchise
After women won the vote in U.S. states, local public-health spending rose at once,
and child mortality from infectious diseases fell by 8--15~\% (Miller 2008).
% ledger: b3pv:mil.suffrage
The inequality prediction has no clean test here. Across 75 countries, average
central-government spending over 2010--2019 \emph{falls} with the average Gini index,
by 0.70 points of GDP per Gini point (s.e.\ 0.13), and by 0.44 (0.14) holding income
per person fixed; but that Gini measures income after taxes and transfers, so a
country that redistributes more shows less inequality: the slope runs at least partly
from spending to the Gini. Across 36 European countries, with the Gini of income
before social transfers, the slope is 0.01 (s.e.\ 0.27): no relation. % figdata: mr_test
Even that Gini is not the model's: it is after taxes, and incomes before transfers
respond to the transfers people expect.
\breaks The poor vote less, and the median voter is richer than the median
household; redistribution also insures against risk, which the model leaves out;
and spending is not all redistribution.
\end{model}

\begin{definition}[Wagner's law; Baumol cost disease]\label{def:b3:public-choice-voting:size-of-government}
\index{Wagner's law}\index{Baumol cost disease}
\emph{Wagner's law} is the measured regularity that public spending rises as a share
of output as income per person grows. \emph{Baumol cost disease} is the rise in the
relative price of services whose productivity grows slowly, such as teaching and
care, when wages rise with productivity elsewhere.
\end{definition}

\begin{proposition}[The relative price of the stagnant sector]\label{prop:b3:public-choice-voting:baumol}
If labour is the only input, wages are equal across two sectors, and output per
worker grows at $g$ in one sector and zero in the other, the relative price of the
stagnant sector's output grows at rate $g$.
\end{proposition}

\begin{proof}
With competitive prices equal to unit labour costs, $p_i = w/a_i$ with $a_i$ output per
worker. The relative price $p_2/p_1 = a_1/a_2$ grows at the growth rate of $a_1$ minus that of $a_2$,
which is $g - 0 = g$.
\end{proof}

Governments buy mostly services --- teaching, care, administration --- so the cost
disease raises the measured share of spending even at constant real provision.

\begin{omchart}{Central-government expense against GDP per person (PPP, constant
dollars, logarithmic scale), 138 countries, 2019.}{World Bank, World Development
Indicators (GC.XPN.TOTL.GD.ZS, NY.GDP.PCAP.PP.KD)}
\begin{tikzpicture}
\begin{semilogxaxis}[omscatter, log ticks with fixed point, width=0.58\linewidth,
  height=0.4\linewidth, xmin=600, xmax=150000, ymin=0, ymax=95,
  xlabel={GDP per person (dollars)}, ylabel={Expense (\% of GDP)}]
  \addplot table[x=y, y=s] {figdata/bachelor-3/21-public-choice-voting/out/wagner.csv};
\end{semilogxaxis}
\end{tikzpicture}
\end{omchart}

% figdata: wagner_stats
Across countries, a doubling of income per person goes with a central-government
share higher by $3.45 \times \ln 2 = 2.4$ points (slope on log income 3.45, robust s.e.\ 0.88), but income
explains only a tenth of the variation ($R^2$ 0.10). Two small island states, Nauru
(90~\% of GDP) and Kiribati (70~\%), pull the slope down: without them it is 4.06.
In short, \omterm{def:b3:public-choice-voting:size-of-government}{Wagner's law} holds on average and
leaves most of the differences to be explained by other things.

\section{Fiscal rules}\label{sec:b3:public-choice-voting:rules}

\begin{definition}[Fiscal rule]\label{def:b3:public-choice-voting:fiscal-rule}
\index{fiscal rule}
A \emph{fiscal rule} is a lasting legal limit on a government's budget aggregates:
a balanced-budget requirement, a cap on debt, spending or taxes.
\end{definition}

\begin{record}[Balanced-budget rules in the United States states]
When regional downturns opened large budget deficits in several states in the late
1980s, states with rules forbidding a deficit to be carried into the next year,
or with tax and spending limits, adjusted faster. States where one party held both
the legislature and the governorship reacted much faster than states under
divided control, and in years when the governor was up for election, tax rises
and spending cuts were significantly smaller. % ledger: b3pv:poterba
\end{record}

\section{Local governments and moving households}\label{sec:b3:public-choice-voting:tiebout}

\begin{definition}[Tiebout sorting]\label{def:b3:public-choice-voting:tiebout}
\index{Tiebout sorting}\index{fiscal competition}
\emph{Tiebout sorting} is the choice by mobile households among many local
jurisdictions according to their bundles of taxes and public services, which sorts
households with similar demands together; the competition among jurisdictions for
residents and firms is \emph{fiscal competition}.
\end{definition}

\begin{model}{Tiebout sorting}
\assumes Households move at low cost among many jurisdictions, each offering a
tax and service bundle.
\predicts Households sort by their demand for local services; differences in
taxes and services are capitalised into house prices.
\tested Comparing houses on opposite sides of school attendance boundaries, which
holds the neighbourhood fixed, Black (1999) found that parents pay 2.5~\% more for a
house whose school's test scores are 5~\% higher: school quality is capitalised, as
the model predicts. % ledger: b3pv:black.boundary
\Cref{exo:b3:public-choice-voting:9} asks what such data would have to show for taxes.
\breaks Zoning limits who can move in, moving is costly, and services such as
schools depend on who the neighbours are, which sorting itself changes.
\end{model}

A household comparing two towns with the same services and a tax difference of $\Delta T$
a year is indifferent when the house in the higher-tax town is cheaper by the
present value of the difference, $\Delta T/r$ at interest rate $r$ for ever: the tax is
capitalised, and the buyer pays it once in the price.

\section{Exercises}\label{sec:b3:public-choice-voting:exercises}

\begin{exercise}[$\star$]\label{exo:b3:public-choice-voting:1}
In the opening chart, compare the rise of the median share from 1870 to 1950 with
the rise from 1950 to 2020. Which country is at the top of the range in 2020?
\end{exercise}

\begin{exercise}[$\star$]\label{exo:b3:public-choice-voting:2}
Five voters' ideal tax rates are 10, 15, 25, 30 and 45~\%. With single-peaked
preferences, which rate wins a majority against every other?
\end{exercise}

\begin{exercise}[$\star$]\label{exo:b3:public-choice-voting:3}
In an election with 1{,}000 other voters, each voting for either candidate with
probability one half, what is the probability of a tie among them?
\end{exercise}

\begin{exercise}[$\star$]\label{exo:b3:public-choice-voting:4}
What would contradict the median-voter prediction in a two-party race?
\end{exercise}

\begin{exercise}[$\star\star$]\label{exo:b3:public-choice-voting:5}
Mean income is 40{,}000 and median 32{,}000; the deadweight parameter $a$ is 0.8. Find the
Meltzer--Richard tax rate. What is it if the median rises to 36{,}000?
\end{exercise}

\begin{exercise}[$\star\star$]\label{exo:b3:public-choice-voting:6}
In the pool example, the \omterm{def:b3:public-choice-voting:agenda-setter}{reversion point} is 1 lane instead of 2. What does the setter
propose? What if the \omterm{def:b3:public-choice-voting:agenda-setter}{reversion point} is 6 lanes?
\end{exercise}

\begin{exercise}[$\star\star$]\label{exo:b3:public-choice-voting:7}
Output per worker grows at 2~\% a year in manufacturing and 0~\% in care. By how much has
the relative price of care risen after 20 years?
\end{exercise}

\begin{exercise}[$\star\star$]\label{exo:b3:public-choice-voting:8}
With the WDI data, estimate the cross-country slope of central-government expense
on log income per person in 2000 and in 2019. Has \omterm{def:b3:public-choice-voting:size-of-government}{Wagner's law} strengthened?
\end{exercise}

\begin{exercise}[$\star\star$]\label{exo:b3:public-choice-voting:9}
What house-price data would contradict the capitalisation of local taxes?
\end{exercise}

\begin{exercise}[$\star\star\star$]\label{exo:b3:public-choice-voting:10}
Prove the Meltzer--Richard comparative static, and show that if the deadweight cost
is $\tfrac a2t^2\bar y$ with $a$ rising in the share of income that can avoid tax, the chosen rate falls as
avoidance grows.
\end{exercise}

\begin{exercise}[$\star\star\star$]\label{exo:b3:public-choice-voting:11}
Three districts vote on two projects. Project A is worth 100 to district 1, project B
100 to district 2; each costs 120, shared equally by the three. Does either pass
alone? Can districts 1 and 2 trade votes to pass both, and is the result efficient?
\end{exercise}

\begin{exercise}[$\star\star\star$]\label{exo:b3:public-choice-voting:12}
With the JST data, regress the central-government spending share on log real GDP
per person with country fixed effects, 1870--2020, and then add year fixed effects.
What survives of \omterm{def:b3:public-choice-voting:size-of-government}{Wagner's law} within countries?
\end{exercise}

\section{Problem: Brenmoor Votes on a Pool}\label{sec:b3:public-choice-voting:problem}

\begin{problem}[{Weekend problem --- Brenmoor votes on a pool}]\label{pb:b3:public-choice-voting:pool}
Brenmoor's households are those of Brenland's survey (\code{load("households")}). A
pool of $q$ lanes costs each household $20q$ dalens a year; a household of $n$ members
values it at $60n\ln q$. The old pool has 2 lanes. % figdata: brenmoor_stats, brenmoor_u

\medskip\noindent\textbf{Part I --- the referendum.}
\begin{enumerate}
\item Find the ideal pool size of a household of $n$ members.
\item Show that preferences are single-peaked in $q$.
\item Find the median household size in the survey.
\item What size wins a referendum against every other?
\item What size would maximise total net benefit, and why does it differ?
\item Who is outvoted, and by how much do they lose?
\end{enumerate}

\medskip\noindent\textbf{Part II --- the agenda setter.} The council alone proposes a
size, and if voters reject it the old pool stays.
\begin{enumerate}[resume]
\item Compute the median household's net benefit at 2 lanes.
\item Find the largest size it accepts.
\item Why does this proposal pass?
\item What would the setter propose if the old pool had 1 lane?
\item What if the old pool had been closed, leaving no pool?
\item What rule would stop the setter exploiting the \omterm{def:b3:public-choice-voting:agenda-setter}{reversion point}?
\end{enumerate}

\medskip\noindent\textbf{Part III --- redistribution.} The council finances a
transfer to every household by a linear income tax, with a deadweight cost
$\tfrac a2t^2\bar y$, $a = 0.5$.
\begin{enumerate}[resume]
\item Find the median and mean gross incomes of the survey.
\item Find the tax rate the median voter chooses.
\item Compute the transfer per household.
\item What happens to the rate if the poorest fifth stop voting?
\item How does the cross-country evidence bear on the model's prediction?
\end{enumerate}

\medskip\noindent\textbf{Part IV --- the next town.} The neighbouring town offers the
same services for a tax 500 dalens a year higher. The interest rate is 4~\%.
\begin{enumerate}[resume]
\item How much cheaper must an identical house there be?
\item Who bears the higher tax: current owners or future buyers?
\item What happens to the difference if households cannot move?
\item How would the referendum on the pool change house prices?
\item How does capitalisation change who has a stake in a vote on local taxes?
\item What evidence would show capitalisation in Brenmoor?
\item State the result: the pool size the referendum chooses, the setter's proposal,
and the capitalisation of the tax difference.
\end{enumerate}
\end{problem}
